% ===================================================================
%  CUADRO N.º 8 — La Cosecha. — La Caza, la Vendimia, la Pesca.
%  Traduction 2026 d'apres texto/io, controlee sur texto/fr et
%  texto/en. Consigne : voir texto/es/10-cuadro-01.tex.
%
%  L'appel de note est dans le TITRE de la scene, comme dans les
%  autres colonnes.
% ===================================================================

%%K t08-noto-1 noto
\VUnotes{143.2mm}{%
\textit{(*) Porque esta palabra no es precisa: ¿se trata de una cosecha
de lino, de una vendimia, de una recolección de manzanas? Quizá fuera
mejor usar «} \VUgras{mesono} \textit{», propuesto en} Mondo \textit{XIV,
p. 242. Pero hasta ahora esa raíz nueva no ha sido adoptada por la
Academia y espera la discusión según las reglas habituales del
movimiento idista.}%
}

%%K t08-apar-1 apar
\VUsaut{5.0mm}
\VUtitre{15.0pt}{60}{CUADRO N.º 8}
\VUsaut{3.0mm}
\VUcentre{13.2pt}{90}{\textit{Primera escena.}}
\VUsaut{3.0mm}
\VUpk{10.2pt}{40}{La Cosecha (*).}
\VUsaut{4.0mm}
\VUpk{10.2pt}{40}{Los Aspectos del Campo.}
\VUsaut{3.0mm}

%%K t08-c1-01-1 p
1. --- Con el \VUgras{verano}, el campo ha cambiado de aspecto una vez más.
La semilla confiada al surco ha germinado; una plantita verde ha salido
de la tierra y ha dado su espiga. El grano se ha formado, luego ha
madurado, y ya no queda más que cosechar.

%%K t08-c1-02-1 p
2. --- Las \VUgras{flores del campo} están abiertas. \VUgras{Acianos}
\textsuperscript{(1)}, \VUgras{amapolas} \textsuperscript{(2)} y \VUgras{dedaleras}
\textsuperscript{(3)} ponen en el tono uniforme de los trigales el alegre
contraste de sus frescas \VUgras{corolas}.

%%K t08-c1-03-1 p
3. --- Los árboles frutales se han cubierto de hojas; el fruto ha
madurado. El pequeño \VUgras{cerezo} \textsuperscript{(4)} está cargado de hermosas
\VUgras{cerezas} \textsuperscript{(5)}, que los niños cogen y comen con deleite.

%%K t08-c1-04-1 p
4. --- El rebaño puede pasar ya todo el día al aire libre.

%%K t08-c1-04-2 p
Los \VUgras{corderos} \textsuperscript{(6)} han crecido y se han convertido en
hermosas \VUgras{ovejas} \textsuperscript{(7)} y hermosos \VUgras{carneros}
\textsuperscript{(8)} de espeso vellón. Pacen tranquilamente la \VUgras{hierba} del
\VUgras{pasto} \textsuperscript{(9)}, salpicado de \VUgras{margaritas}.

%%K t08-c1-04-3 p
Pronto empezarán a esquilar a estos animales para tener su \VUgras{lana},
con la que se fabrica el paño de nuestra ropa.

%%K t08-tit-2 sub
\VUsaut{5.0mm}
\VUpk{10.2pt}{40}{El Pastor.}
\VUsaut{3.4mm}

%%K t08-c1-05-1 p
5. --- El \VUgras{pastor} \textsuperscript{(10)} no parece hacerles mucho caso. Solo
se preocupa de la tormenta que pasa por el horizonte, y el
\VUgras{zagal} \textsuperscript{(13)}, que se lleva la \VUgras{cantimplora}
\textsuperscript{(14)} a los labios, parece aún más despreocupado.

%%K t08-c1-06-1 p
6. --- El calor es agobiante; porque el \VUgras{perro} \textsuperscript{(15)} viene
también a refrescarse; lame el agua fresca del \VUgras{arroyo}
\textsuperscript{(16)} y espanta a una pobre \VUgras{rana} \textsuperscript{(17)} que estaba
agazapada en la hierba de la orilla. Esta salta al agua clara donde
florecen los \VUgras{nenúfares} \textsuperscript{(18)}.

%%K t08-c1-07-1 p
7. --- Una \VUgras{aguzanieves} \textsuperscript{(19)} se baña junto al pequeño
\VUgras{manantial} \textsuperscript{(20)} que brota entre dos rocas y se esconde bajo
la \VUgras{maleza espinosa} \textsuperscript{(21)} y los \VUgras{matorrales de
espino albar} \textsuperscript{(22)}.

%%K t08-tit-3 sub
\VUsaut{5.0mm}
\VUpk{10.2pt}{40}{La Cosecha.}
\VUsaut{3.4mm}

%%K t08-c1-08-1 p
8. --- Para los niños del campo, la siega del trigo es una época muy
divertida. Dan alegres \VUgras{volteretas} \textsuperscript{(24)} sobre los
\VUgras{montones de paja} \textsuperscript{(23)}, ayudan a los \VUgras{segadores}
\textsuperscript{(26)} y, mientras estos siegan el \VUgras{trigal} \textsuperscript{(27)} con
sus \VUgras{guadañas} \textsuperscript{(25)}, bien afiladas en la \VUgras{piedra de
afilar} \textsuperscript{(30)}, ellos recogen el trigo y lo atan en \VUgras{gavillas}
\textsuperscript{(31)} o en \VUgras{haces} \textsuperscript{(32)}.

%%K t08-c1-09-1 p
9. --- A veces se permite a los mayores cortar con una \VUgras{hoz}
\textsuperscript{(12)} los \VUgras{tallos dorados} \textsuperscript{(28)} que llevan las
pesadas \VUgras{espigas} \textsuperscript{(29)} llenas de grano; y los pequeños,
mientras guardan el \VUgras{ganado} \textsuperscript{(33)} que pace en el
\VUgras{prado} \textsuperscript{(34)} a la sombra del gran \VUgras{tilo} \textsuperscript{(35)},
cazan las hermosas \VUgras{mariposas} con su \VUgras{red} \textsuperscript{(36)}.

%%K t08-c1-09-2 p
Algunos hasta se suben al \VUgras{carro} \textsuperscript{(37)} para colocar las
gavillas que les pasan con una \VUgras{horca} \textsuperscript{(38)}.

%%K t08-c1-10-1 p
10. --- Por toda la \VUgras{llanura} \textsuperscript{(39)}, donde alzan el vuelo las
\VUgras{alondras} \textsuperscript{(41)}, reina la animación. Todos están ocupados
con la siega. Pero, mientras los pobres siguen usando las herramientas
antiguas --- \VUgras{guadañas, hoces} y \VUgras{mayales} \textsuperscript{(40)} ---,
los propietarios ricos hacen segar su trigo o su \VUgras{centeno}
\textsuperscript{(43)} con una \VUgras{segadora mecánica} \textsuperscript{(42)}. Luego, sin
necesidad de llevar las gavillas al granero, se hacen allí mismo grandes
\VUgras{almiares} \textsuperscript{(44)}.

%%K t08-c1-11-1 p
11. --- Entonces se trae una \VUgras{trilladora} \textsuperscript{(47)} que una
\VUgras{locomóvil} \textsuperscript{(45)} mueve por medio de una \VUgras{correa de
transmisión} \textsuperscript{(46)}, y así el grano queda separado de la
\VUgras{paja} sin salir del campo en que fue sembrado.

%%K t08-tit-4 sub
\VUsaut{5.0mm}
\VUpk{10.2pt}{40}{La Tormenta.}
\VUsaut{3.4mm}

%%K t08-c1-12-1 p
12. --- A menudo estallan violentas \VUgras{tormentas} \textsuperscript{(53)} en la
época de la siega. Primero se oyen a lo lejos los truenos del
\VUgras{trueno}; se levanta un fuerte \VUgras{viento del oeste}
\textsuperscript{(54)} que dobla con quejidos los árboles más gruesos del
\VUgras{bosque} \textsuperscript{(49)}. Negros \VUgras{nubarrones} \textsuperscript{(56)}
asaltan el cielo; los atraviesan los zigzags cegadores del
\VUgras{relámpago} \textsuperscript{(55)}. La \VUgras{lluvia} y a veces el
\VUgras{granizo} empiezan a caer, aplastando la cosecha lista para ser
segada.

%%K t08-c1-13-1 p
13. --- A menudo el rayo incendia algún edificio. Así, el año pasado el
tejado del \VUgras{molino de viento} \textsuperscript{(50)} quedó reducido a cenizas
y dos de sus \VUgras{aspas} \textsuperscript{(51)} quedaron destrozadas.

%%K t08-c1-14-1 p
14. --- Al cabo de unas horas vuelve la calma. Un brillante \VUgras{arco
iris} \textsuperscript{(52)} aparece en el horizonte y la naturaleza reanuda su
vida, interrumpida un instante. Los campesinos, que se habían refugiado
en las \VUgras{cabañas} \textsuperscript{(48)}, vuelven al trabajo y se esfuerzan
valientemente en reparar los daños que acaban de sufrir.

%%K t08-tit-5 sub
\VUsaut{6.0mm}
\VUcentre{13.2pt}{90}{\textit{Segunda escena.}}
\VUsaut{3.6mm}

%%K t08-tit-6 sub
\VUpk{10.2pt}{40}{La Caza. --- La Vendimia.}
\VUsaut{1.4mm}
\VUpk{10.2pt}{40}{La Pesca.}
\VUsaut{4.0mm}

%%K t08-c2-01-1 p
1. --- Hacia finales de \VUgras{agosto} o a mediados de
\VUgras{septiembre}, según las regiones, se abre la veda. Apenas han
hecho salir los primeros rayos del sol a las \VUgras{lagartijas}
\textsuperscript{(2)} de su agujero, cuando el \VUgras{cazador} \textsuperscript{(5)} se
pone en camino con sus \VUgras{perros} \textsuperscript{(4)}.

%%K t08-c2-02-1 p
2. --- Se ha puesto su sólido \VUgras{traje de caza} \textsuperscript{(9)} de paño
grueso y unas \VUgras{polainas} de cuero con las que podrá andar por la
hierba mojada donde se esconden los \VUgras{sapos} \textsuperscript{(1)}; ha cogido
su \VUgras{escopeta} \textsuperscript{(6)} y su \VUgras{morral} \textsuperscript{(10)}, y ya se
ha marchado.

%%K t08-c2-03-1 p
3. --- El ardor de la persecución lo lleva a un viñedo donde la vendimia
está en pleno apogeo. Los hijos del viñador, que suelen guardar las
cabras en el camino, hoy las dejan pacer a su antojo y, después de atar a
una estaca a un viejo \VUgras{macho cabrío} \textsuperscript{(3)} que no dejaría de
aprovechar su libertad para escaparse, se dan también su parte de placer.
Uno coge un racimo de uvas de una \VUgras{cesta de vendimia} \textsuperscript{(12)} y
se divierte haciendo correr el \VUgras{zumo} \textsuperscript{(14)} en un \VUgras{vaso}
\textsuperscript{(13)} para hacer mosto (vino sin fermentar). El otro se corona con
una \VUgras{guirnalda de hojas} \textsuperscript{(15)} que acaba de trenzar. Junto a
ellos, unos \VUgras{gorriones} \textsuperscript{(16)} descarados llegan hasta el
mismo lagar a picotear las uvas.

%%K t08-c2-04-1 p
4. --- Mientras los niños se divierten, los hombres trabajan. Los
\VUgras{vendimiadores} \textsuperscript{(18)} llevan la uva a la bodega en
\VUgras{cuévanos} \textsuperscript{(17)}; un \VUgras{tonelero} \textsuperscript{(19)} arregla
las \VUgras{barricas} \textsuperscript{(20)}, comprueba si las \VUgras{duelas}
\textsuperscript{(22)} y los \VUgras{fondos} \textsuperscript{(21)} están en buen estado y
sustituye los \VUgras{aros} \textsuperscript{(23)} rotos. También fue él quien hizo
las grandes \VUgras{cubas} \textsuperscript{(25)} en las que se echa la \VUgras{uva}
\textsuperscript{(26)} para que fermente.

%%K t08-c2-05-1 p
5. --- Cuando la fermentación ha terminado, se saca el vino y se pone el
resto en el \VUgras{lagar} \textsuperscript{(28)} y, apretando poco a poco el gran
\VUgras{husillo} \textsuperscript{(29)}, se exprime el zumo, que corre a un
\VUgras{cubeto} \textsuperscript{(32)}, y se obtiene todavía más \VUgras{vino}
\textsuperscript{(31)}.

%%K t08-tit-8 sub
\VUsaut{6.0mm}
\VUpk{10.2pt}{40}{Cerveza y Sidra.}
\VUsaut{3.4mm}

%%K t08-c2-06-1 p
6. --- Por la pared de la \VUgras{bodega} \textsuperscript{(27)} trepa una rama de
\VUgras{lúpulo} \textsuperscript{(30)}. Aquí no es más que una planta de adorno,
pero en algunos países, donde la vid es muy rara, el lúpulo se cultiva
para fabricar la \VUgras{cerveza}.

%%K t08-c2-07-1 p
7. --- El pequeño \VUgras{manzano} \textsuperscript{(33)} está cargado de manzanas
que, cuando estén maduras, darán una \VUgras{sidra} excelente. En su
\VUgras{jaula} \textsuperscript{(34)}, el \VUgras{canario} \textsuperscript{(35)} y el
\VUgras{jilguero} \textsuperscript{(36)} miran con envidia a los gorriones que
retozan en libertad.

%%K t08-c2-08-1 p
8. --- Hasta donde alcanza la vista, en la ladera de las colinas, se
alinean los \VUgras{rodrigones} \textsuperscript{(40)} que sostienen las magníficas
\VUgras{cepas} \textsuperscript{(39)}, cargadas de pesados \VUgras{racimos}.

%%K t08-c2-08-2 p
Durante todo el año, el \VUgras{viñador} \textsuperscript{(42)} ha trabajado para
cuidar su \VUgras{viñedo} \textsuperscript{(38)}, y ahora se ve recompensado de su
esfuerzo con una hermosa cosecha. Las \VUgras{vendimiadoras}
\textsuperscript{(41)} llenan las cubas puestas en el carro que tiran las
\VUgras{mulas} \textsuperscript{(43)}, y todos están alegres.

%%K t08-tit-9 sub
\VUsaut{5.0mm}
\VUpk{10.2pt}{40}{La Pesca.}
\VUsaut{3.4mm}

%%K t08-c2-09-1 p
9. --- Lo mismo ocurre con los \VUgras{pescadores de caña} \textsuperscript{(44)},
que desde por la mañana han venido a instalarse a la orilla del agua con
sus \VUgras{cañas de pescar} \textsuperscript{(45)}.

%%K t08-c2-09-2 p
El \VUgras{pez} \textsuperscript{(47)} pica en el anzuelo en cuanto echan su
\VUgras{sedal} \textsuperscript{(46)} al agua, y apenas les da tiempo a renovar el
\VUgras{cebo} cuando una nueva víctima colea en la punta del hilo.

%%K t08-c2-09-3 p
Con todo, el \VUgras{pescador de red} \textsuperscript{(48)}, que desde su
\VUgras{barca} \textsuperscript{(49)} echa la esparavel en medio del \VUgras{río}
\textsuperscript{(50)}, coge muchos más.

%%K t08-c2-10-1 p
10. --- El otoño es la estación de los buenos paseos: a pie, a
\VUgras{caballo} \textsuperscript{(52)}, en \VUgras{bicicleta} \textsuperscript{(53)}, en
\VUgras{automóvil} \textsuperscript{(54)}. Los \VUgras{caminos} \textsuperscript{(51)} están
limpios y se aprovechan los últimos días buenos, porque ya las
\VUgras{golondrinas} \textsuperscript{(56)} se reúnen en los tejados para volar a
todo ala hacia países más cálidos. Las hojas del viejo \VUgras{nogal}
\textsuperscript{(57)} amarillean, las \VUgras{nueces} caen y los altos
\VUgras{árboles} agrupados en \VUgras{bosquecillo} \textsuperscript{(58)} sobre la
\VUgras{altura} \textsuperscript{(59)} pronto se quedarán sin hojas.

%%K t08-c2-11-1 p
11. --- El \VUgras{sol} \textsuperscript{(60)} sale más tarde, los días menguan,
empieza a sentirse un frío bastante vivo por la \VUgras{mañana}, y ya
bandadas de \VUgras{becadas} \textsuperscript{(61)} dejan nuestros climas
inhóspitos.
\VUsaut{6.0mm}
\VUfilet{22mm}
